\chapter{Chapter07 — CSP Algorithms}
%%% generated by the blueprint pipeline from FABL.Chapter07.CSPAlgorithms
\section{Overview}

Declarations of the Lean module \texttt{FABL.Chapter07.CSPAlgorithms}.

\section{Declarations}

\begin{definition}[Exact-one promise problem]
\lean{FABL.EncodedMaximizationProblem.exactOnePromise}
\statementsource{Analysis-of-Boolean-Functions-by-Ryan-ODonnell}{pdf-ffbe9c7b3306-p188-b004}
\label{FABL.EncodedMaximizationProblem.exactOnePromise}
\leanok
\uses{FABL.BinaryPromiseProblem, FABL.EncodedMaximizationProblem}
For an encoded maximization problem, the yes-instances are the well-formed inputs whose
stated optimum is exactly $1$, and the no-instances are the well-formed inputs whose
stated optimum is less than $1$. Inputs outside these two sets are not assigned either
promise.
\end{definition}

\begin{definition}[Search by One-One Restriction Queries]
\lean{FABL.OneOneRestrictionSearch}
\label{FABL.OneOneRestrictionSearch}
\leanok
\uses{FABL.EncodedMaximizationProblem, FABL.EncodedMaximizationProblem.exactOnePromise, FABL.PolytimeMap, FABL.PolytimeMap.DistinguishesPromise}
Let an encoded maximization problem consist of a language of valid Boolean-string
instances, a feasibility relation between instances and Boolean-string solutions, a
rational objective value for each instance-solution pair, and a rational optimum for
each instance; assume every feasible solution on a valid instance has value at most that
optimum. A one-one restriction search for this problem is a polynomial-time map from
decision procedures to reconstruction procedures such that, for every decision procedure
and valid instance, the reconstructed solution is feasible, and every valid instance has
optimum at most $1$. Moreover, whenever the decision procedure distinguishes the
problem’s exact-one promise, every valid instance whose optimum is $1$ has a
reconstructed solution of value exactly $1$.
\end{definition}

\begin{theorem}[Approximation from an Exact-One Distinguisher]
\lean{FABL.OneOneRestrictionSearch.approximationOfDistinguisher}
\label{FABL.OneOneRestrictionSearch.approximationOfDistinguisher}
\leanok
\uses{FABL.EncodedMaximizationProblem, FABL.EncodedMaximizationProblem.exactOnePromise, FABL.OneOneRestrictionSearch, FABL.PolytimeMap, FABL.PolytimeMap.DistinguishesPromise, FABL.PolytimeMap.IsThresholdApproximation}
\difficulty{3}
Let an encoded maximization problem admit a one-one restriction search, and let a
polynomial-time map distinguish its exact-one promise, outputting true on every valid
instance with optimum $1$ and false on every valid instance with optimum less than $1$.
Then the solution reconstructed from that map is feasible on every valid instance;
moreover, it has value at least $1$ whenever the instance's optimum is at least $1$.
\end{theorem}

\begin{definition}[Maximum-cut assignment]
\lean{FABL.maxCutSatisfyingAssignment}
\statementsource{vershynin-hdp-book}{
  pdf-681b6f3d947f-p073-b004,
  pdf-681b6f3d947f-p073-b005,
  pdf-681b6f3d947f-p074-b004
}
\statementsource{Analysis-of-Boolean-Functions-by-Ryan-ODonnell}{
  pdf-ffbe9c7b3306-p184-b003,
  pdf-ffbe9c7b3306-p185-b004,
  pdf-ffbe9c7b3306-p186-b001
}
\label{FABL.maxCutSatisfyingAssignment}
\leanok
\uses{FABL.CSPInstance, FABL.CSPInstance.optimalAssignment, FABL.Sign, FABL.maxCutTemplate}
Let $V$ be a finite set and let $P$ be an instance of the maximum-cut constraint
template on $V$. The maximum-cut satisfying assignment associated with $P$ is the
assignment from $V$ to the signs $\{ -1,1\}$ specified as $P$'s optimal assignment.
\end{definition}

\begin{theorem}[Satisfiable cut instances are fully satisfied]
\lean{FABL.maxCutSatisfyingAssignment_fullySatisfied}
\label{FABL.maxCutSatisfyingAssignment_fullySatisfied}
\leanok
\uses{FABL.CSPInstance, FABL.CSPInstance.FullySatisfied, FABL.CSPInstance.Satisfiable, FABL.CSPInstance.value_eq_one_iff_fullySatisfied, FABL.maxCutSatisfyingAssignment, FABL.maxCutTemplate}
\difficulty{2}
Let $V$ be a finite set, and let $P$ be a maximum-cut constraint instance on $V$ with a
nonempty multiset of constraints in which every variable occurs in the scope of some
constraint. If the optimum value of $P$ is $1$, then an optimal assignment satisfies
every constraint of $P$.
\end{theorem}

\begin{definition}[Encoded maximum-cut problem]
\lean{FABL.encodedMaxCutProblem}
\statementsource{vershynin-hdp-book}{
  pdf-681b6f3d947f-p073-b004,
  pdf-681b6f3d947f-p073-b005,
  pdf-681b6f3d947f-p074-b003,
  pdf-681b6f3d947f-p074-b004
}
\statementsource{Analysis-of-Boolean-Functions-by-Ryan-ODonnell}{
  pdf-ffbe9c7b3306-p184-b003,
  pdf-ffbe9c7b3306-p185-b004,
  pdf-ffbe9c7b3306-p185-b005,
  pdf-ffbe9c7b3306-p185-b006,
  pdf-ffbe9c7b3306-p186-b001
}
\label{FABL.encodedMaxCutProblem}
\leanok
\uses{FABL.BinaryEncoding, FABL.BinaryEncoding.language, FABL.CSPInstance, FABL.CSPInstance.optimum, FABL.CSPInstance.value, FABL.CSPInstance.value_le_optimum, FABL.EncodedMaximizationProblem, FABL.Sign, FABL.maxCutTemplate}
Let $V$ be finite, and choose canonical binary encodings for maximum-cut instances on
$V$ and assignments of signs to $V$. The well-formed inputs are the canonical encodings
of instances. An input and output are feasible exactly when both decode successfully;
their value is the instance value of the decoded assignment when both decode, and is $0$
otherwise. The stated optimum is the decoded instance's optimum when the input decodes,
and is $0$ otherwise.
\end{definition}

\begin{definition}[Polynomial-time maximum-cut search]
\lean{FABL.MaxCutPolytimeSearch}
\statementsource{vershynin-hdp-book}{
  pdf-681b6f3d947f-p074-b005,
  pdf-681b6f3d947f-p074-b006,
  pdf-681b6f3d947f-p075-b001
}
\statementsource{Analysis-of-Boolean-Functions-by-Ryan-ODonnell}{pdf-ffbe9c7b3306-p189-b001}
\label{FABL.MaxCutPolytimeSearch}
\leanok
\uses{FABL.BinaryEncoding, FABL.CSPInstance, FABL.PolytimeMap, FABL.Sign, FABL.maxCutSatisfyingAssignment, FABL.maxCutTemplate}
Let $V$ be finite, with canonical binary encodings for maximum-cut instances on $V$ and
assignments of signs to $V$. A polynomial-time maximum-cut search structure consists of
a polynomial-time map on binary strings that, on the encoding of every instance $P$,
returns the encoding of $P$'s specified maximum-cut satisfying assignment.
\end{definition}

\begin{theorem}[Is threshold approximation]
\lean{FABL.MaxCutPolytimeSearch.isThresholdApproximation}
\statementsource{Analysis-of-Boolean-Functions-by-Ryan-ODonnell}{
  pdf-ffbe9c7b3306-p188-b004,
  pdf-ffbe9c7b3306-p189-b001
}
\label{FABL.MaxCutPolytimeSearch.isThresholdApproximation}
\leanok
\uses{FABL.BinaryEncoding, FABL.CSPInstance, FABL.CSPInstance.Satisfiable, FABL.CSPInstance.optimum_le_one, FABL.CSPInstance.value_eq_one_iff_fullySatisfied, FABL.MaxCutPolytimeSearch, FABL.PolytimeMap.IsThresholdApproximation, FABL.Sign, FABL.encodedMaxCutProblem, FABL.maxCutSatisfyingAssignment, FABL.maxCutSatisfyingAssignment_fullySatisfied, FABL.maxCutTemplate}
\difficulty{5}
% dataset-pipeline: stub -- prose pending build_informal_statements.py
Statement pending regeneration.
\end{theorem}

\begin{definition}[Finite randomized assignment]
\lean{FABL.FiniteRandomizedAssignment}
\statementsource{Analysis-of-Boolean-Functions-by-Ryan-ODonnell}{pdf-ffbe9c7b3306-p198-b006}
\label{FABL.FiniteRandomizedAssignment}
\leanok
Let $V$ and $D$ be finite sets. A finite randomized assignment on $V$ with values in $D$
specifies a probability law on $D$ for each variable in $V$.
\end{definition}

\begin{definition}[Joint Law of a Random Assignment]
\lean{FABL.FiniteRandomizedAssignment.sampleLaw}
\label{FABL.FiniteRandomizedAssignment.sampleLaw}
\leanok
\uses{FABL.FiniteRandomizedAssignment, FABL.independentProductPMF}
Let $V$ and $D$ be finite sets, and assign to each variable $v\in V$ a probability law
on $D$. The sampling law on assignments $V\to D$ is the joint law in which the variables
are independent and each variable has its assigned marginal law.
\end{definition}

\begin{definition}[Expected payoff]
\lean{FABL.FiniteRandomizedAssignment.expectedPayoff}
\statementsource{Analysis-of-Boolean-Functions-by-Ryan-ODonnell}{pdf-ffbe9c7b3306-p198-b006}
\label{FABL.FiniteRandomizedAssignment.expectedPayoff}
\leanok
\uses{FABL.FiniteRandomizedAssignment, FABL.FiniteRandomizedAssignment.sampleLaw, FABL.pmfExpectation}
Let $V$ and $D$ be finite sets, and let a randomized assignment specify a marginal
probability law on $D$ for each variable in $V$. For a payoff function from assignments
$V\to D$ to $\mathbb{R}$, its expected payoff is the expectation under the assignment's
sample law.
\end{definition}

\begin{theorem}[Expected Payoff Identity]
\lean{FABL.FiniteRandomizedAssignment.expectedPayoff_sample}
\label{FABL.FiniteRandomizedAssignment.expectedPayoff_sample}
\leanok
\uses{FABL.FiniteRandomizedAssignment, FABL.FiniteRandomizedAssignment.expectedPayoff, FABL.FiniteRandomizedAssignment.sampleLaw, FABL.pmfExpectation}
\difficulty{1}
Let $V$ and $D$ be finite sets, let each variable $v\in V$ have a probability law on
$D$, and let $f:(V\to D)\to\mathbb{R}$ be a payoff. The expectation of $f$ under the
independent product law of these variable laws equals the assignment’s expected payoff.
\end{theorem}

\begin{definition}[Ordered Product of Probability Laws]
\lean{FABL.orderedProductPMF}
\label{FABL.orderedProductPMF}
\leanok
Let $D$ be a finite set, let $n$ be a nonnegative integer, and assign a probability law
on $D$ to each coordinate in $\mathrm{Fin}(n)$. The ordered product law on coordinate
assignments is formed by sampling the coordinates independently according to their
respective laws; when $n=0$, it is concentrated on the unique assignment on the empty
coordinate set.
\end{definition}

\begin{definition}[A Maximizer on a Finite Space]
\lean{FABL.finiteMaximizer}
\label{FABL.finiteMaximizer}
\leanok
Let $\Omega$ be a finite nonempty set and let $f:\Omega\to\mathbb R$. A finite maximizer
of $f$ is an element of $\Omega$ at which $f$ attains its maximum.
\end{definition}

\begin{theorem}[Maximum on a finite set]
\lean{FABL.le_finiteMaximizer}
\label{FABL.le_finiteMaximizer}
\leanok
\uses{FABL.finiteMaximizer}
\difficulty{1}
Let $\Omega$ be a finite nonempty set, let $f:\Omega\to\mathbb R$, and let $x\in\Omega$.
Then $f(x)$ is at most the value of $f$ at a maximizer of $f$ on $\Omega$.
\end{theorem}

\begin{theorem}[Expectation bounded by a maximizing value]
\lean{FABL.pmfExpectation_le_finiteMaximizer}
\label{FABL.pmfExpectation_le_finiteMaximizer}
\leanok
\uses{FABL.finiteMaximizer, FABL.le_finiteMaximizer, FABL.pmfExpectation, FABL.pmfExpectation_const, FABL.pmfExpectation_mono}
\difficulty{3}
Let $\Omega$ be a finite nonempty set, let $p$ be a probability mass function on
$\Omega$, and let $f:\Omega\to\mathbb{R}$. The expectation of $f$ under $p$ is at most
the value of $f$ at a maximizer of $f$ on $\Omega$.
\end{theorem}

\begin{definition}[Conditional-expectation assignment]
\lean{FABL.conditionalExpectationAssignment}
\statementsource{Analysis-of-Boolean-Functions-by-Ryan-ODonnell}{pdf-ffbe9c7b3306-p199-b005}
\label{FABL.conditionalExpectationAssignment}
\leanok
\uses{FABL.finiteMaximizer, FABL.orderedProductPMF, FABL.pmfExpectation}
Let $D$ be a finite nonempty set, let $n$ be a nonnegative integer, assign a probability
law on $D$ to each of the $n$ coordinates, and let $f:(\mathrm{Fin}(n)\to D)\to\mathbb{R}$ be a payoff. The recursively selected assignment fixes each coordinate to a value maximizing
the expected payoff over the remaining coordinates under their independent product law,
conditional on the coordinates already fixed; when $n=0$, it is the unique empty
assignment.
\end{definition}

\begin{theorem}[Conditional assignment dominates expected payoff]
\lean{FABL.pmfExpectation_le_conditionalExpectationAssignment}
\proofsource{Analysis-of-Boolean-Functions-by-Ryan-ODonnell}{pdf-ffbe9c7b3306-p199-b005}
\label{FABL.pmfExpectation_le_conditionalExpectationAssignment}
\leanok
\uses{FABL.conditionalExpectationAssignment, FABL.finiteMaximizer, FABL.orderedProductPMF, FABL.pmfExpectation, FABL.pmfExpectation_bind, FABL.pmfExpectation_const, FABL.pmfExpectation_le_finiteMaximizer, FABL.pmfExpectation_map}
\difficulty{5}
Let $D$ be a finite nonempty set, let $n\geq 0$, and assign a probability mass function
$p_i$ on $D$ to each coordinate $i\in\{0,\ldots,n-1\}$. For every payoff function
$F:D^n\to\mathbb{R}$, its expectation under independent coordinate sampling is at most
its value at the assignment chosen by successively maximizing the conditional expected
payoff of the remaining coordinates.
\end{theorem}

\begin{theorem}[Independent product of uniform distributions]
\lean{FABL.independentProductPMF_uniform}
\label{FABL.independentProductPMF_uniform}
\leanok
\uses{FABL.independentProductPMF, FABL.uniformPMF}
\difficulty{3}
Let $V$ and $D$ be finite sets, with $D$ nonempty. The independent product of uniform
distributions on $D$, one for each coordinate in $V$, is the uniform distribution on the
set of functions from $V$ to $D$.
\end{theorem}

\begin{definition}[Uniform randomized assignment]
\lean{FABL.uniformRandomizedAssignment}
\statementsource{Analysis-of-Boolean-Functions-by-Ryan-ODonnell}{pdf-ffbe9c7b3306-p199-b002}
\label{FABL.uniformRandomizedAssignment}
\leanok
\uses{FABL.FiniteRandomizedAssignment, FABL.uniformPMF}
For finite sets $V$ and $D$, with $D$ nonempty, the uniform randomized assignment on $V$
with values in $D$ assigns to each variable the uniform probability distribution on $D$.
\end{definition}

\begin{theorem}[Sample law of a uniform randomized assignment]
\lean{FABL.uniformRandomizedAssignment_sampleLaw}
\label{FABL.uniformRandomizedAssignment_sampleLaw}
\leanok
\uses{FABL.FiniteRandomizedAssignment.sampleLaw, FABL.independentProductPMF_uniform, FABL.uniformPMF, FABL.uniformRandomizedAssignment}
\difficulty{1}
Let $V$ and $D$ be finite sets, with $D$ nonempty. If each variable in $V$ is
independently assigned a uniform value from $D$, then the resulting probability law on
assignments $V\to D$ is the uniform distribution on the set of all such assignments.
\end{theorem}

\begin{definition}[Splitting Assignments Along a Constraint Scope]
\lean{FABL.cspScopeSplitEquiv}
\label{FABL.cspScopeSplitEquiv}
\leanok
\uses{FABL.CSPConstraint, FABL.CSPTemplate}
Let $V$ and $D$ be sets, let $\Psi$ be a constraint template on $D$, and let $c$ be a
constraint on $V$ with predicate $p$ and an injective scope map from the predicate’s
coordinate set into $V$. Every assignment $V\to D$ is equivalently specified by its
values on the ordered variables in the scope and its values on the variables outside the
scope; thus there is a bijection between assignments $V\to D$ and pairs consisting of an
assignment to the coordinates of $p$ and an assignment to the variables outside the
scope.
\end{definition}

\begin{theorem}[Constraint Scope Restriction]
\lean{FABL.cspScopeSplitEquiv_fst_apply}
\label{FABL.cspScopeSplitEquiv_fst_apply}
\leanok
\uses{FABL.CSPConstraint, FABL.CSPTemplate, FABL.cspScopeSplitEquiv}
\difficulty{3}
Let $D$ be a domain, let $V$ be a finite set of variables, and let a constraint have an
ordered scope of pairwise distinct variables. For any assignment $a:V\to D$, the first
component of its decomposition into values on the scope and values off the scope, at
position $i$ in the constraint's scope, is $a$ evaluated at the variable occupying
position $i$.
\end{theorem}

\begin{definition}[Triples as Functions on Three Coordinates]
\lean{FABL.finThreeArrowEquiv}
\label{FABL.finThreeArrowEquiv}
\leanok
For any set $D$, the set of functions from $\mathrm{Fin}(3)$ to $D$ is in bijection with
$D\times D\times D$, by sending a function to its values at coordinates $0$, $1$, and
$2$. The inverse assigns those three entries to the corresponding coordinates.
\end{definition}

\begin{theorem}[Encoding false in the binary field]
\lean{FABL.boolF₂Equiv_false}
\leanok
\uses{FABL.boolF₂Equiv}
\difficulty{1}
The encoding of the Boolean value false in $\mathbb{F}_2$ is $0$.
\end{theorem}

\begin{theorem}[Encoding true in the binary field]
\lean{FABL.boolF₂Equiv_true}
\leanok
\uses{FABL.boolF₂Equiv}
\difficulty{1}
The encoding of the Boolean value true in $\mathbb{F}_2$ is $1$.
\end{theorem}

\begin{theorem}[Acceptance probability of a three-variable linear predicate]
\lean{FABL.expect_e3LinPredicate}
\label{FABL.expect_e3LinPredicate}
\leanok
\uses{FABL.boolF₂Equiv, FABL.boolF₂Equiv_false, FABL.boolF₂Equiv_true, FABL.boolVectorF₂CubeEquiv, FABL.finThreeArrowEquiv, FABL.maxE3LinTemplate, FABL.𝔽₂}
\difficulty{5}
For any right-hand side $r\in\mathbb{F}_2$, a uniformly random triple of elements of
$\mathbb{F}_2$ satisfies the predicate given by the maximum three-variable linear
template with right-hand side $r$ with probability $1/2$.
\end{theorem}

\begin{theorem}[Expected satisfaction of a random linear constraint]
\lean{FABL.expect_satisfiedIndicator_maxE3Lin}
\label{FABL.expect_satisfiedIndicator_maxE3Lin}
\leanok
\uses{FABL.CSPConstraint, FABL.CSPConstraint.satisfiedIndicator, FABL.cspScopeSplitEquiv, FABL.cspScopeSplitEquiv_fst_apply, FABL.expect_e3LinPredicate, FABL.maxE3LinTemplate, FABL.𝔽₂}
\difficulty{4}
Let $V$ be a finite set and consider a constraint on three distinct variables in $V$
requiring their sum in $\mathbb{F}_2$ to equal a specified right-hand side in
$\mathbb{F}_2$. Under a uniformly random assignment of elements of $\mathbb{F}_2$ to
$V$, the expected satisfaction indicator of this constraint is $1/2$.
\end{theorem}

\begin{theorem}[Expected number of satisfied linear constraints]
\lean{FABL.expect_satisfiedIndicators_maxE3Lin}
\label{FABL.expect_satisfiedIndicators_maxE3Lin}
\leanok
\uses{FABL.CSPConstraint, FABL.CSPConstraint.satisfiedIndicator, FABL.expect_satisfiedIndicator_maxE3Lin, FABL.maxE3LinTemplate, FABL.𝔽₂}
\difficulty{4}
Let $V$ be a finite set and let $\mathcal{C}$ be a finite multiset of constraints, each
requiring the sum in $\mathbb{F}_2$ of three distinct variables to equal a specified
right-hand side in $\mathbb{F}_2$. Under a uniformly random assignment of elements of
$\mathbb{F}_2$ to $V$, the expected number of satisfied constraints in $\mathcal{C}$ is
$|\mathcal{C}|/2$.
\end{theorem}

\begin{theorem}[Expected value of a three-variable linear system]
\lean{FABL.expect_value_maxE3Lin}
\label{FABL.expect_value_maxE3Lin}
\leanok
\uses{FABL.CSPInstance, FABL.CSPInstance.satisfiedCount, FABL.CSPInstance.value, FABL.expect_satisfiedIndicators_maxE3Lin, FABL.maxE3LinTemplate, FABL.𝔽₂}
\difficulty{3}
Let $V$ be a finite set and consider a nonempty multiset of constraints, each specifying
three distinct variables of $V$ and a right-hand side in the field with two elements; a
constraint is satisfied by an assignment $a:V\to\mathbb{F}_2$ exactly when the sum of
its three variable values equals its right-hand side. Suppose every variable in $V$
occurs in the scope of at least one constraint. The expected fraction of satisfied
constraints, counting multiplicities, under a uniformly random assignment
$a:V\to\mathbb{F}_2$ is $1/2$.
\end{theorem}

\begin{theorem}[Expected value of a random 3-SAT assignment]
\lean{FABL.uniformRandomizedAssignment_expectedValue_maxE3Sat}
\label{FABL.uniformRandomizedAssignment_expectedValue_maxE3Sat}
\leanok
\uses{FABL.CSPInstance, FABL.CSPInstance.value, FABL.FiniteRandomizedAssignment.sampleLaw, FABL.expect_value_maxE3Sat, FABL.maxE3SatTemplate, FABL.pmfExpectation, FABL.pmfExpectation_uniformPMF_eq_expect, FABL.uniformRandomizedAssignment, FABL.uniformRandomizedAssignment_sampleLaw}
\difficulty{3}
Let $V$ be a finite set of Boolean variables, and let a 3-SAT instance have a nonempty
multiset of clauses, with every variable occurring in some clause. If each variable is
assigned an independent uniform Boolean value, then the expected fraction of satisfied
clauses, counting multiplicities, is $7/8$.
\end{theorem}

\begin{theorem}[Uniform randomized assignment expected value max e3 lin]
\lean{FABL.uniformRandomizedAssignment_expectedValue_maxE3Lin}
\statementsource{Analysis-of-Boolean-Functions-by-Ryan-ODonnell}{pdf-ffbe9c7b3306-p199-b002}
\label{FABL.uniformRandomizedAssignment_expectedValue_maxE3Lin}
\leanok
\uses{FABL.CSPInstance, FABL.CSPInstance.value, FABL.FiniteRandomizedAssignment.sampleLaw, FABL.expect_value_maxE3Lin, FABL.maxE3LinTemplate, FABL.pmfExpectation, FABL.pmfExpectation_uniformPMF_eq_expect, FABL.uniformRandomizedAssignment, FABL.uniformRandomizedAssignment_sampleLaw, FABL.𝔽₂}
\difficulty{3}
% dataset-pipeline: stub -- prose pending build_informal_statements.py
Statement pending regeneration.
\end{theorem}

\begin{definition}[Even-Parity Boolean Triples]
\lean{FABL.evenBooleanTriple}
\label{FABL.evenBooleanTriple}
\leanok
For each of the four elements of $\mathrm{Fin}(4)$, the corresponding Boolean triple is,
in order, $(0,0,0)$, $(0,1,1)$, $(1,0,1)$, or $(1,1,0)$. These are exactly the Boolean
triples with an even number of true coordinates.
\end{definition}

\begin{definition}[Negation bit for a linear clause]
\lean{FABL.e3LinClauseNegation}
\label{FABL.e3LinClauseNegation}
\leanok
\uses{FABL.boolF₂Equiv, FABL.evenBooleanTriple, FABL.𝔽₂}
For a right-hand side $r\in\mathbb F_2$, a clause label $c\in\{0,1,2,3\}$, and a
position $i\in\{0,1,2\}$, define the negation bit to be the exclusive-or of the
condition that $r=0$ and $i=0$ with the bit indicating whether the $i$th coordinate of
the even Boolean triple indexed by $c$ is true.
\end{definition}

\begin{theorem}[Count of satisfying replacement clauses]
\lean{FABL.sum_e3LinReplacementClauses}
\proofsource{Analysis-of-Boolean-Functions-by-Ryan-ODonnell}{pdf-ffbe9c7b3306-p200-b008}
\label{FABL.sum_e3LinReplacementClauses}
\leanok
\uses{FABL.boolF₂Equiv, FABL.e3LinClauseNegation, FABL.e3SatPredicate, FABL.evenBooleanTriple, FABL.maxE3LinTemplate, FABL.𝔽₂}
\difficulty{5}
Let $r\in\mathbb{F}_2$ and let $(v_0,v_1,v_2)\in\{0,1\}^3$. Form four clauses, one for
each even-parity triple $e\in\{0,1\}^3$, by negating the first literal of $e$ exactly
when $r=0$ and leaving its other literals as in $e$. A clause is satisfied when at least
one of its three literals evaluates to true. The number of these four clauses satisfied
by $(v_0,v_1,v_2)$ is $3$ plus $1$ if $v_0+v_1+v_2=r$ in $\mathbb{F}_2$, and is $3$
otherwise.
\end{theorem}

\begin{definition}[Clause from a linear constraint]
\lean{FABL.e3LinConstraintToE3Sat}
\statementsource{Analysis-of-Boolean-Functions-by-Ryan-ODonnell}{pdf-ffbe9c7b3306-p200-b008}
\label{FABL.e3LinConstraintToE3Sat}
\leanok
\uses{FABL.CSPConstraint, FABL.e3LinClauseNegation, FABL.maxE3LinTemplate, FABL.maxE3SatTemplate}
Given a constraint over $\mathbb{F}_2$ with three distinct variables and a right-hand
side, and an index $j\in\{0,1,2,3\}$, form a three-literal Boolean clause on the same
ordered variables. Its literal signs are those specified by the $j$th clause in the
four-clause replacement of the linear constraint.
\end{definition}

\begin{definition}[Four clauses from a linear constraint]
\lean{FABL.e3LinConstraintClauses}
\statementsource{Analysis-of-Boolean-Functions-by-Ryan-ODonnell}{pdf-ffbe9c7b3306-p200-b008}
\label{FABL.e3LinConstraintClauses}
\leanok
\uses{FABL.CSPConstraint, FABL.e3LinConstraintToE3Sat, FABL.maxE3LinTemplate, FABL.maxE3SatTemplate}
Given an E3-Lin constraint over $\mathbb{F}_2$ on three distinct variables, associate
the multiset of four E3-Sat clauses obtained by applying the four-clause replacement,
one clause for each of the four indices in $\{0,1,2,3\}$.
\end{definition}

\begin{theorem}[Clause count for a linear constraint]
\lean{FABL.card_e3LinConstraintClauses}
\label{FABL.card_e3LinConstraintClauses}
\leanok
\uses{FABL.CSPConstraint, FABL.e3LinConstraintClauses, FABL.maxE3LinTemplate}
\difficulty{1}
For any constraint requiring three distinct variables over the two-element field, the
multiset of three-literal E3-SAT constraints constructed from it has cardinality $4$.
\end{theorem}

\begin{theorem}[Satisfied clauses in the linear equation encoding]
\lean{FABL.sum_satisfiedIndicator_e3LinConstraintClauses}
\proofsource{Analysis-of-Boolean-Functions-by-Ryan-ODonnell}{pdf-ffbe9c7b3306-p200-b008}
\label{FABL.sum_satisfiedIndicator_e3LinConstraintClauses}
\leanok
\uses{FABL.CSPConstraint, FABL.CSPConstraint.satisfiedIndicator, FABL.boolF₂Equiv, FABL.e3LinClauseNegation, FABL.e3LinConstraintClauses, FABL.e3SatPredicate, FABL.maxE3LinTemplate, FABL.sum_e3LinReplacementClauses}
\difficulty{1}
Let $V$ be a set, let a constraint specify a right-hand side in $\mathbb F_2$ and an
ordered triple of distinct variables from $V$, and let $a:V\to\{0,1\}$ be any Boolean
assignment. Form the four three-literal clauses associated with that linear constraint.
The number of those clauses satisfied by $a$ is three plus the indicator that the
corresponding equation, that the sum of the three assigned values equals the specified
right-hand side in $\mathbb F_2$, is satisfied.
\end{theorem}

\begin{definition}[Clause replacement for a constraint multiset]
\lean{FABL.e3LinConstraintsToE3Sat}
\statementsource{Analysis-of-Boolean-Functions-by-Ryan-ODonnell}{pdf-ffbe9c7b3306-p200-b008}
\label{FABL.e3LinConstraintsToE3Sat}
\leanok
\uses{FABL.CSPConstraint, FABL.e3LinConstraintClauses, FABL.maxE3LinTemplate, FABL.maxE3SatTemplate}
Given a multiset of E3-Lin constraints over $\mathbb{F}_2$, replace each constraint by
its multiset of four E3-Sat clauses, preserving multiplicities, and take the resulting
multiset of clauses.
\end{definition}

\begin{theorem}[Size of the Converted Constraint Multiset]
\lean{FABL.card_e3LinConstraintsToE3Sat}
\label{FABL.card_e3LinConstraintsToE3Sat}
\leanok
\uses{FABL.CSPConstraint, FABL.e3LinConstraintsToE3Sat, FABL.maxE3LinTemplate}
\difficulty{4}
For any multiset of three-variable linear constraints over $\mathbb{F}_2$, converting
each constraint into its four three-variable E3-SAT clauses yields a multiset whose
cardinality is exactly four times the cardinality of the original multiset.
\end{theorem}

\begin{definition}[E3-Lin to E3-Sat instance reduction]
\lean{FABL.e3LinToE3Sat}
\statementsource{Analysis-of-Boolean-Functions-by-Ryan-ODonnell}{
  pdf-ffbe9c7b3306-p200-b008,
  pdf-ffbe9c7b3306-p201-b001
}
\label{FABL.e3LinToE3Sat}
\leanok
\uses{FABL.CSPInstance, FABL.card_e3LinConstraintsToE3Sat, FABL.e3LinConstraintClauses, FABL.e3LinConstraintToE3Sat, FABL.e3LinConstraintsToE3Sat, FABL.maxE3LinTemplate, FABL.maxE3SatTemplate}
Every E3-Lin instance over $\mathbb{F}_2$ is sent to an E3-Sat instance on the same
variables by replacing each linear constraint with its four-clause multiset. The
resulting instance has at least one constraint, and every declared variable occurs in
the scope of some constraint.
\end{definition}

\begin{theorem}[Satisfied-count identity for the constraint reduction]
\lean{FABL.satisfiedCount_e3LinConstraintsToE3Sat}
\proofsource{Analysis-of-Boolean-Functions-by-Ryan-ODonnell}{pdf-ffbe9c7b3306-p200-b008}
\label{FABL.satisfiedCount_e3LinConstraintsToE3Sat}
\leanok
\uses{FABL.CSPConstraint, FABL.CSPConstraint.satisfiedIndicator, FABL.boolF₂Equiv, FABL.e3LinConstraintsToE3Sat, FABL.maxE3LinTemplate, FABL.sum_satisfiedIndicator_e3LinConstraintClauses}
\difficulty{4}
Let $V$ be a type, let a finite multiset of three-variable linear equations over the
field with two elements be given, and let $a:V\to\{\mathrm{false},\mathrm{true}\}$ be a
Boolean assignment. The total number of satisfied three-literal clauses produced by
replacing each equation with its four associated clauses equals three times the number
of original equations plus the number of original equations satisfied by the assignment
mapped to the field with two elements.
\end{theorem}

\begin{definition}[{Boolean assignment as an $$F₂-valued function}]
\lean{FABL.boolAssignmentToF₂}
\leanok
\uses{FABL.boolF₂Equiv, FABL.𝔽₂}
Given a set $V$ of variables and a Boolean assignment
$a:V\to\{\mathrm{false},\mathrm{true}\}$, define its corresponding $\mathbb F_2$-valued
assignment by applying the standard identification of Booleans with $\mathbb F_2$ at
each variable.
\end{definition}

\begin{definition}[F₂-valued assignment as a Boolean function]
\lean{FABL.f₂AssignmentToBool}
\leanok
\uses{FABL.boolF₂Equiv, FABL.𝔽₂}
Given a set $V$ of variables and an assignment $a:V\to\mathbb F_2$, define the
corresponding Boolean assignment by applying the inverse of the standard identification
of Booleans with $\mathbb F_2$ at each variable.
\end{definition}

\begin{theorem}[Boolean and binary field assignments]
\lean{FABL.boolAssignmentToF₂_f₂AssignmentToBool}
\leanok
\uses{FABL.boolAssignmentToF₂, FABL.boolF₂Equiv, FABL.f₂AssignmentToBool, FABL.𝔽₂}
\difficulty{3}
Let $V$ be any set and let $a:V\to\mathbb{F}_2$, where $\mathbb{F}_2$ is the field with
two elements. Converting each value of $a$ to a Boolean and then converting back to
$\mathbb{F}_2$ recovers $a$.
\end{theorem}

\begin{theorem}[Boolean assignment round trip]
\lean{FABL.f₂AssignmentToBool_boolAssignmentToF₂}
\leanok
\uses{FABL.boolAssignmentToF₂, FABL.boolF₂Equiv, FABL.f₂AssignmentToBool}
\difficulty{3}
For any set $V$ and any assignment of Boolean values to the elements of $V$, converting
the assignment to values in $\mathbb{F}_2$ and then back to Boolean values recovers the
original assignment.
\end{theorem}

\begin{theorem}[Value under parity-to-satisfaction conversion]
\lean{FABL.value_e3LinToE3Sat}
\proofsource{Analysis-of-Boolean-Functions-by-Ryan-ODonnell}{
  pdf-ffbe9c7b3306-p200-b008,
  pdf-ffbe9c7b3306-p201-b001
}
\label{FABL.value_e3LinToE3Sat}
\leanok
\uses{FABL.CSPInstance, FABL.CSPInstance.satisfiedCount, FABL.CSPInstance.value, FABL.boolAssignmentToF₂, FABL.boolF₂Equiv, FABL.card_e3LinConstraintsToE3Sat, FABL.e3LinToE3Sat, FABL.maxE3LinTemplate, FABL.satisfiedCount_e3LinConstraintsToE3Sat}
\difficulty{4}
Let $P$ be a nonempty finite collection of three-variable linear equations over $\mathbb
F_2$, with every declared variable appearing in at least one equation, and let $x$ be a
Boolean assignment to its variables. Convert each equation into the four three-literal
clauses specified by the even Boolean triples, with the clause signs determined by the
equation's right-hand side, and let $P'$ be the resulting instance. The fraction of
clauses in $P'$ satisfied by $x$ equals $\frac34$ plus one quarter of the fraction of
equations in $P$ satisfied by the corresponding assignment over $\mathbb F_2$.
\end{theorem}

\begin{theorem}[Optimum e3 lin to e3 sat]
\lean{FABL.optimum_e3LinToE3Sat}
\statementsource{Analysis-of-Boolean-Functions-by-Ryan-ODonnell}{pdf-ffbe9c7b3306-p201-b001}
\label{FABL.optimum_e3LinToE3Sat}
\leanok
\uses{FABL.CSPInstance, FABL.CSPInstance.exists_value_eq_optimum, FABL.CSPInstance.optimum, FABL.CSPInstance.value, FABL.CSPInstance.value_le_optimum, FABL.boolAssignmentToF₂, FABL.boolAssignmentToF₂_f₂AssignmentToBool, FABL.e3LinToE3Sat, FABL.f₂AssignmentToBool, FABL.maxE3LinTemplate, FABL.value_e3LinToE3Sat}
\difficulty{4}
% dataset-pipeline: stub -- prose pending build_informal_statements.py
Statement pending regeneration.
\end{theorem}

\begin{theorem}[Value transfer from satisfiability to linear constraints]
\lean{FABL.half_add_delta_le_e3LinValue_of_sevenEighths_add_delta_le}
\label{FABL.half_add_delta_le_e3LinValue_of_sevenEighths_add_delta_le}
\leanok
\uses{FABL.CSPInstance, FABL.CSPInstance.value, FABL.boolAssignmentToF₂, FABL.e3LinToE3Sat, FABL.maxE3LinTemplate, FABL.value_e3LinToE3Sat}
\difficulty{3}
Let $P$ be an instance of the three-variable linear constraint problem over
$\bbr/	ext{?}$—that is, over the field with two elements—and let $a$ be a Boolean
assignment to its variables. If $\delta\footnote{}$ is a nonnegative rational number and
the value of the associated three-literal satisfiability instance under $a$ is at least
$7/8+\delta$, then the value of $P$ under the corresponding assignment over the field
with two elements is at least $1/2+\delta$.
\end{theorem}

\begin{theorem}[One sub delta le e3 sat optimum of e3 lin]
\lean{FABL.one_sub_delta_le_e3SatOptimum_of_e3Lin}
\statementsource{Analysis-of-Boolean-Functions-by-Ryan-ODonnell}{pdf-ffbe9c7b3306-p201-b001}
\label{FABL.one_sub_delta_le_e3SatOptimum_of_e3Lin}
\leanok
\uses{FABL.CSPInstance, FABL.CSPInstance.optimum, FABL.e3LinToE3Sat, FABL.maxE3LinTemplate, FABL.optimum_e3LinToE3Sat}
\difficulty{3}
% dataset-pipeline: stub -- prose pending build_informal_statements.py
Statement pending regeneration.
\end{theorem}

\begin{definition}[Algorithm induced by an E3-Sat algorithm]
\lean{FABL.e3LinAlgorithmOfE3Sat}
\statementsource{Analysis-of-Boolean-Functions-by-Ryan-ODonnell}{pdf-ffbe9c7b3306-p201-b001}
\label{FABL.e3LinAlgorithmOfE3Sat}
\leanok
\uses{FABL.CSPInstance, FABL.boolAssignmentToF₂, FABL.e3LinToE3Sat, FABL.maxE3LinTemplate, FABL.maxE3SatTemplate, FABL.𝔽₂}
Given an algorithm that assigns a Boolean value to each variable of every E3-Sat
instance, define an algorithm for E3-Lin instances by applying it to the instance
produced by the four-clause replacement and interpreting each Boolean output as an
element of $\mathbb{F}_2$.
\end{definition}

\begin{theorem}[Approximation transfer from three-literal SAT]
\lean{FABL.e3LinAlgorithmOfE3Sat_isApproximation}
\proofsource{Analysis-of-Boolean-Functions-by-Ryan-ODonnell}{pdf-ffbe9c7b3306-p201-b001}
\label{FABL.e3LinAlgorithmOfE3Sat_isApproximation}
\leanok
\uses{FABL.CSPInstance, FABL.CSPInstance.optimum, FABL.CSPInstance.value, FABL.e3LinAlgorithmOfE3Sat, FABL.e3LinToE3Sat, FABL.half_add_delta_le_e3LinValue_of_sevenEighths_add_delta_le, FABL.maxE3LinTemplate, FABL.maxE3SatTemplate, FABL.one_sub_delta_le_e3SatOptimum_of_e3Lin}
\difficulty{4}
Let $V$ be a finite set of variables, let $\delta\in\mathbb{Q}$ satisfy $0\leq\delta$,
and let $A$ be an algorithm that assigns Boolean values to variables for every
three-literal satisfiability instance over $V$. Suppose that on every such instance $Q$,
if its optimum satisfied fraction is at least $1-\delta$, then $A(Q)$ satisfies a
fraction at least $7/8+\delta$. Then, for every three-variable linear constraint
instance $P$ over $\mathbb{F}_2$ on $V$, if its optimum satisfied fraction is at least
$1-\delta$, the assignment produced by applying $A$ to the associated three-literal
instance and interpreting its Boolean values in $\mathbb{F}_2$ satisfies a fraction at
least $1/2+\delta$ of the constraints of $P$.
\end{theorem}

\begin{definition}[Polynomial-time reduction adapter]
\lean{FABL.E3LinE3SatPolytimeAdapter}
\statementsource{Analysis-of-Boolean-Functions-by-Ryan-ODonnell}{pdf-ffbe9c7b3306-p201-b001}
\label{FABL.E3LinE3SatPolytimeAdapter}
\leanok
\uses{FABL.EncodedMaximizationProblem, FABL.PolytimeMap}
Let two encoded maximization problems and a rational number $\delta$ be given. An
adapter consists of polynomial-time maps reducing instances and recovering assignments,
such that well-formed E3-Lin instances map to well-formed E3-Sat instances; feasible
E3-Sat solutions recover to feasible E3-Lin solutions; an E3-Lin optimum of at least
$1-\delta$ implies an E3-Sat optimum of at least $1-\delta$; and any E3-Sat solution of
value at least $7/8+\delta$ recovers to an E3-Lin solution of value at least
$1/2+\delta$.
\end{definition}

\begin{theorem}[Approximation transfers through the adapter]
\lean{FABL.E3LinE3SatPolytimeAdapter.transferApproximation}
\proofsource{Analysis-of-Boolean-Functions-by-Ryan-ODonnell}{pdf-ffbe9c7b3306-p201-b001}
\label{FABL.E3LinE3SatPolytimeAdapter.transferApproximation}
\leanok
\uses{FABL.E3LinE3SatPolytimeAdapter, FABL.EncodedMaximizationProblem, FABL.PolytimeMap, FABL.PolytimeMap.IsThresholdApproximation, FABL.PolytimeMap.comp}
\difficulty{4}
Let $L$ and $S$ be encoded maximization problems and let $\delta$ be a rational number.
Suppose an adapter reduces instances of $L$ to instances of $S$, recovers an
$L$-feasible solution from every feasible solution to the reduced instance, transfers
the condition $1-\delta \leq \operatorname{OPT}_L(x)$ to
$1-\delta \leq \operatorname{OPT}_S(R(x))$, and transfers value at least
$7/8+\delta$ on a reduced solution to value at least $1/2+\delta$ on its
recovered solution. If a polynomial-time algorithm for $S$ always outputs a feasible
solution and, whenever the optimum is at least $1-\delta$, achieves value at least
$7/8+\delta$, then the algorithm obtained by reducing an $L$ instance, applying that
algorithm, and recovering its solution always outputs an $L$-feasible solution and
achieves value at least $1/2+\delta$ whenever the $L$ optimum is at least $1-\delta$.
\end{theorem}

\begin{definition}[Even sign triples]
\lean{FABL.evenSignTriple}
\label{FABL.evenSignTriple}
\leanok
\uses{FABL.Sign}
For each label $c\in\{0,1,2,3\}$, define a triple of signs by the four possibilities
$(1,1,1)$, $(1,-1,-1)$, $(-1,1,-1)$, and $(-1,-1,1)$, respectively.
\end{definition}

\begin{definition}[Falsifying sign triple]
\lean{FABL.e3LinFalsifyingSign}
\label{FABL.e3LinFalsifyingSign}
\leanok
\uses{FABL.Sign, FABL.evenSignTriple}
For a sign $b\in\{-1,1\}$ and a label $c\in\{0,1,2,3\}$, define a triple of signs whose
first coordinate is $-b$ times the first coordinate of the even sign triple indexed by
$c$, and whose other two coordinates agree with that even sign triple.
\end{definition}

\begin{definition}[Sign clause excluding one triple]
\lean{FABL.e3SatSignClause}
\label{FABL.e3SatSignClause}
\leanok
\uses{FABL.Sign, FABL.SignPredicate}
Given a falsifying triple of signs, define a three-input Boolean predicate that accepts
precisely those triples of signs whose value differs from the falsifying triple in at
least one coordinate.
\end{definition}

\begin{definition}[Predicates from the E3-Sat sign clauses]
\lean{FABL.maxE3SatSignPredicates}
\statementsource{Analysis-of-Boolean-Functions-by-Ryan-ODonnell}{pdf-ffbe9c7b3306-p185-b003}
\label{FABL.maxE3SatSignPredicates}
\leanok
\uses{FABL.SignPredicate, FABL.e3SatSignClause}
The set of sign predicates for E3-Sat is the range of the construction that associates a
sign predicate to each E3-Sat clause.
\end{definition}

\begin{theorem}[Finiteness of signed E3-SAT predicates]
\lean{FABL.maxE3SatSignPredicates_finite}
\label{FABL.maxE3SatSignPredicates_finite}
\leanok
\uses{FABL.e3SatSignClause, FABL.maxE3SatSignPredicates}
\difficulty{1}
For each finite arity bound, the family of maximum E3-SAT sign predicates is finite.
\end{theorem}

\begin{theorem}[Count of satisfying sign replacement clauses]
\lean{FABL.sum_e3SatSignReplacementClauses}
\label{FABL.sum_e3SatSignReplacementClauses}
\leanok
\uses{FABL.Sign, FABL.e3LinFalsifyingSign, FABL.e3LinPredicate, FABL.e3SatSignClause, FABL.evenSignTriple}
\difficulty{5}
Let $b\in\{-1,1\}$ and let $(v_0,v_1,v_2)\in\{-1,1\}^3$. Form four clauses indexed by
the four even-parity sign triples $e\in\{-1,1\}^3$, whose coordinate product is $+1$. In
the clause indexed by $e$, the falsifying triple is $e$ with its first coordinate
multiplied by $-b$; the clause is satisfied unless the input triple equals that
falsifying triple. The number of these clauses satisfied by $(v_0,v_1,v_2)$ is $3$ plus
$1$ if $v_0v_1v_2=b$, and is $3$ otherwise.
\end{theorem}

\begin{theorem}[Average randomized clause value]
\lean{FABL.expect_e3SatSignClause_randomizedValue}
\label{FABL.expect_e3SatSignClause_randomizedValue}
\leanok
\uses{FABL.Sign, FABL.SignPredicate.randomizedValue, FABL.e3LinFalsifyingSign, FABL.e3LinPredicate, FABL.e3SatSignClause, FABL.fintypeExpectation_pmfExpectation_comm, FABL.independentProductPMF, FABL.pmfExpectation, FABL.pmfExpectation_add, FABL.pmfExpectation_const, FABL.pmfExpectation_const_mul, FABL.signPMFOfMean, FABL.sum_e3SatSignReplacementClauses}
\difficulty{5}
Let $b\in\{\pm1\}$ and let $\mu_0,\mu_1,\mu_2\in[-1,1]$. Average over the four triples
of signs whose product is $1$ the randomized acceptance value of the clause that accepts
unless its three inputs equal the corresponding falsifying signs (with the first
falsifying sign adjusted by $b$). This average equals $3/4$ plus one quarter of the
randomized acceptance value of the predicate that accepts when the product of three
independently randomized signs equals $b$, where the signs have respective means
$\mu_0,\mu_1,\mu_2$.
\end{theorem}

\begin{definition}[Three-query linear sign tester]
\lean{FABL.E3LinSignTester}
\statementsource{Analysis-of-Boolean-Functions-by-Ryan-ODonnell}{
  pdf-ffbe9c7b3306-p191-b002,
  pdf-ffbe9c7b3306-p192-b003
}
\label{FABL.E3LinSignTester}
\leanok
\uses{FABL.Sign}
For a natural number $n$, an $n$-dimensional three-query linear sign tester consists of
a finite seed set with a probability distribution, a sign specifying the right-hand side
for each seed, and three queries in $\{−1,1\}^n$ selected for each seed.
\end{definition}

\begin{definition}[Conversion to a sign tester]
\lean{FABL.E3LinSignTester.toSignTester}
\label{FABL.E3LinSignTester.toSignTester}
\leanok
\uses{FABL.E3LinSignTester, FABL.SignTester, FABL.e3LinPredicate}
Given an $E3$-linear sign tester with finite seed set, seed distribution, a sign
right-hand side for each seed, and three queries for each seed, define a sign tester
with the same seed set and distribution. Its check at a seed uses the predicate
determined by that seed's right-hand side and the same three queries.
\end{definition}

\begin{definition}[Conversion to a three-query satisfiability tester]
\lean{FABL.E3LinSignTester.toE3SatTester}
\statementsource{Analysis-of-Boolean-Functions-by-Ryan-ODonnell}{pdf-ffbe9c7b3306-p201-b002}
\label{FABL.E3LinSignTester.toE3SatTester}
\leanok
\uses{FABL.E3LinSignTester, FABL.SignTester, FABL.e3LinFalsifyingSign, FABL.e3SatSignClause, FABL.uniformPMF}
Given an $n$-dimensional three-query linear sign tester, form a tester whose random seed
consists of a seed of the original tester together with a uniformly chosen element of a
four-element set. For each such pair, its check uses the original three queries and the
three-query satisfiability clause determined by the original seed’s right-hand-side sign
and the chosen element.
\end{definition}

\begin{theorem}[Converted tester uses three-literal clauses]
\lean{FABL.E3LinSignTester.toE3SatTester_usesMaxE3Sat}
\label{FABL.E3LinSignTester.toE3SatTester_usesMaxE3Sat}
\leanok
\uses{FABL.E3LinSignTester, FABL.E3LinSignTester.toE3SatTester, FABL.SignTester.UsesPredicates, FABL.e3LinFalsifyingSign, FABL.maxE3SatSignPredicates}
\difficulty{3}
For every natural number $n$ and every E3-Lin sign tester on the $n$-dimensional Boolean
cube, each predicate used by its converted E3-Sat tester is a three-literal clause: it
accepts exactly when at least one of the three queried signs differs from its
corresponding falsifying sign.
\end{theorem}

\begin{theorem}[Acceptance probability under tester conversion]
\lean{FABL.E3LinSignTester.acceptanceProbability_toE3SatTester}
\proofsource{Analysis-of-Boolean-Functions-by-Ryan-ODonnell}{pdf-ffbe9c7b3306-p201-b002}
\label{FABL.E3LinSignTester.acceptanceProbability_toE3SatTester}
\leanok
\uses{FABL.BoundedCubeFunction, FABL.E3LinSignTester, FABL.E3LinSignTester.toE3SatTester, FABL.E3LinSignTester.toSignTester, FABL.SignPredicate.randomizedValue, FABL.SignQueryCheck.randomizedValue, FABL.SignTester.acceptanceProbability, FABL.e3LinFalsifyingSign, FABL.e3LinPredicate, FABL.e3SatSignClause, FABL.expect_e3SatSignClause_randomizedValue, FABL.pmfExpectation, FABL.pmfExpectation_add, FABL.pmfExpectation_bind, FABL.pmfExpectation_const, FABL.pmfExpectation_const_mul, FABL.pmfExpectation_map, FABL.pmfExpectation_uniformPMF_eq_expect, FABL.uniformPMF}
\difficulty{5}
For every natural number $n$, every $n$-variable E3-Lin sign tester $T$, and every
function $f$ on the $n$-dimensional Boolean cube with values in $[-1,1]$, the acceptance
probability of the converted E3-Sat tester is $\frac{3}{4}+\frac{1}{4}$ times the
acceptance probability of the converted sign tester.
\end{theorem}

\begin{definition}[Scaled replacement error]
\lean{FABL.e3SatReplacementError}
\label{FABL.e3SatReplacementError}
\leanok
For a function $e:\mathbb R\to\mathbb R$ and a real number $\varepsilon$, define the
replacement error to be $e(\varepsilon)/4$.
\end{definition}

\begin{theorem}[Affine conversion to a satisfiability test]
\lean{FABL.e3LinDictatorTest_to_e3Sat_affine}
\label{FABL.e3LinDictatorTest_to_e3Sat_affine}
\leanok
\uses{FABL.E3LinSignTester, FABL.E3LinSignTester.acceptanceProbability_toE3SatTester, FABL.E3LinSignTester.toE3SatTester, FABL.E3LinSignTester.toE3SatTester_usesMaxE3Sat, FABL.E3LinSignTester.toSignTester, FABL.IsDictatorVsNoNotablesTest, FABL.e3SatReplacementError, FABL.maxE3LinPredicates, FABL.maxE3SatSignPredicates, FABL.maxE3SatSignPredicates_finite}
\difficulty{5}
Let $T$ be a family of three-variable linear-equation testers, one for each positive
number of variables, and suppose the corresponding sign testers form a
dictator-versus-no-notables test with parameters $\alpha,\beta$ and error function
$\lambda$. Then the associated three-literal satisfiability testers form such a test for
the family of all three-literal disjunctions, with lower parameter $3/4+\alpha/4$, upper
parameter $3/4+\beta/4$, and error function $\lambda(\varepsilon)/4$.
\end{theorem}

\begin{theorem}[Conversion to a three-literal satisfiability test]
\lean{FABL.e3LinDictatorTest_to_e3Sat}
\proofsource{Analysis-of-Boolean-Functions-by-Ryan-ODonnell}{
  pdf-ffbe9c7b3306-p201-b002,
  pdf-ffbe9c7b3306-p201-b003
}
\label{FABL.e3LinDictatorTest_to_e3Sat}
\leanok
\uses{FABL.E3LinSignTester, FABL.E3LinSignTester.toE3SatTester, FABL.E3LinSignTester.toSignTester, FABL.IsDictatorVsNoNotablesTest, FABL.e3LinDictatorTest_to_e3Sat_affine, FABL.e3SatReplacementError, FABL.maxE3LinPredicates, FABL.maxE3SatSignPredicates}
\difficulty{4}
Let $T$ be a family of three-variable linear-equation testers, one for each positive
number of variables, and suppose the corresponding sign testers form a
dictator-versus-no-notables test with parameters $0<\alpha<\beta\leq 1$ and error
function $\lambda$. If $3/4+\alpha/4<\beta$, then replacing each tester by its
associated three-literal satisfiability tester gives a dictator-versus-no-notables test
with allowed predicates consisting of all three-literal disjunctions, lower parameter
$3/4+\alpha/4$, upper parameter $\beta$, and error function $\lambda(\varepsilon)/4$.
\end{theorem}
